\documentclass[a4paper,11pt]{article}

\usepackage[utf8]{inputenc}
\usepackage[T1]{fontenc}
\usepackage[french]{babel}
\usepackage{graphicx}
\usepackage{geometry}
\usepackage{hyperref}
\usepackage{enumitem}
\geometry{margin=2.5cm}

% Place tes captures dans le dossier docs/images/ puis decommente les \includegraphics.

\title{\textbf{Rapport de projet --- Application Web IRVE}}
\author{Trinôme 9}
\date{Projet Big Data / IA / Web --- Partie Développement Web --- FISE3 2026}

\begin{document}

\maketitle
\tableofcontents
\newpage

% ======================================================
\section{Contexte et objectif}

Ce projet s'inscrit dans le cadre du projet transversal Big Data / Intelligence
Artificielle / Développement Web. L'objectif de la partie web est de concevoir et
développer une application complète permettant d'étudier les données de la base
nationale des Infrastructures de Recharge pour Véhicules Électriques (IRVE).

L'application permet de visualiser les points de charge et les stations, de consulter
des statistiques par département, et d'exploiter les modèles d'intelligence artificielle
développés lors des phases précédentes du projet (clustering géographique et
classification).

% ======================================================
\section{Technologies utilisées}

Conformément au cahier des charges, nous avons utilisé exclusivement les technologies
imposées :

\begin{itemize}[noitemsep]
  \item \textbf{Front-end} : HTML, CSS et JavaScript (vanilla), avec les bibliothèques
        autorisées Leaflet (cartographie) et Plotly (graphiques).
  \item \textbf{Back-end} : PHP, communiquant avec une base de données MySQL/MariaDB.
  \item \textbf{Échanges client-serveur} : requêtes AJAX (fetch) et données au format JSON.
  \item \textbf{Intégration IA} : appel des scripts Python depuis PHP pour les prédictions.
\end{itemize}

L'application respecte une mise en page commune (en-tête et pied de page identiques sur
toutes les pages), une architecture de code séparée par fonctionnalité, et un code
commenté et réexploitable.

% ======================================================
\section{Architecture de la base de données}

La base de données \texttt{irve\_web} est composée de \textbf{5 tables} :

\begin{itemize}[noitemsep]
  \item \textbf{station} : entité centrale, contenant les informations de chaque station
        (nom, adresse, coordonnées, commune, etc.).
  \item \textbf{point\_de\_charge} : les bornes individuelles rattachées à une station
        (puissance, types de prises, modes de paiement) ainsi que les champs remplis par
        l'IA (cluster, prédictions).
  \item \textbf{amenageur} : l'entité qui installe la station.
  \item \textbf{operateur} : l'entité qui exploite la station.
  \item \textbf{departement} : table de référence pour le regroupement géographique.
\end{itemize}

Les relations sont les suivantes :
\begin{itemize}[noitemsep]
  \item une station est \emph{aménagée} par un aménageur (1,1) et \emph{exploitée} par un
        opérateur (1,1) ;
  \item une station est \emph{située} dans un département (1,1) ;
  \item une station \emph{possède} plusieurs points de charge (1,n).
\end{itemize}

\begin{figure}[h]
  \centering
  % \includegraphics[width=0.9\textwidth]{images/mcd.png}
  \caption{Modèle Conceptuel de Données (MCD)}
\end{figure}

L'import des données depuis le fichier \texttt{export\_IA.csv} est réalisé par un script
Python (\texttt{import\_csv.py}) qui déduplique les aménageurs, opérateurs et départements,
et insère un échantillon des données conformément au cahier des charges.

% ======================================================
\section{Présentation des fonctionnalités}

\subsection{Page d'accueil}

La page d'accueil présente le projet avec un menu de navigation, un descriptif rapide et
une image illustrative. L'en-tête et le pied de page y sont définis et réutilisés sur
toutes les pages.

\begin{figure}[h]
  \centering
  % \includegraphics[width=0.9\textwidth]{images/accueil.png}
  \caption{Page d'accueil de l'application}
\end{figure}

\subsection{Visualisation des points de charge et des stations}

Cette page affiche l'ensemble des points de charge dans un tableau (avec un défilement
interne pour l'ergonomie) et les stations sur une carte interactive Leaflet. Le tableau
présente 7 informations par point de charge (identifiant, station, puissance, types de
prises, gratuité, paiement). Les marqueurs de la carte affichent au clic le nom de la
station, son adresse, son type d'implantation et son nombre de bornes.

\begin{figure}[h]
  \centering
  % \includegraphics[width=0.9\textwidth]{images/tableau.png}
  \caption{Tableau des points de charge}
\end{figure}

\begin{figure}[h]
  \centering
  % \includegraphics[width=0.9\textwidth]{images/carte.png}
  \caption{Carte des stations}
\end{figure}

\subsection{Statistiques par département}

Cette page permet de sélectionner un département dans un menu déroulant et d'afficher
quatre statistiques : le nombre de stations, le nombre de points de charge, la puissance
nominale moyenne, et la répartition des stations par type d'implantation présentée sous
forme de camembert (Plotly).

\begin{figure}[h]
  \centering
  % \includegraphics[width=0.9\textwidth]{images/statistiques.png}
  \caption{Statistiques par département}
\end{figure}

\subsection{Prédiction du cluster des points de charge}

Un bouton « Prédire les clusters », situé sous le tableau des points de charge, ouvre une
nouvelle page. Celle-ci appelle, côté serveur, le script Python de clustering qui prédit
le cluster géographique de chaque point de charge. Les résultats sont enregistrés en base
et les points sont affichés sur une carte, colorés selon leur cluster d'appartenance. Le
clic sur un point affiche son identifiant, sa station, sa puissance et son cluster.

Pour des raisons de performance, le script Python a été adapté en version « batch » afin
de prédire l'ensemble des points en un seul appel.

\begin{figure}[h]
  \centering
  % \includegraphics[width=0.9\textwidth]{images/clusters.png}
  \caption{Carte des points de charge colorés par cluster}
\end{figure}

% ======================================================
\section{Organisation du code}

Le code est structuré par fonctionnalité dans les dossiers suivants :

\begin{itemize}[noitemsep]
  \item \texttt{config/} : connexion à la base de données (PDO) ;
  \item \texttt{php/api/} : les interfaces client-serveur (un fichier PHP par requête) ;
  \item \texttt{assets/css/} : la feuille de style commune ;
  \item \texttt{assets/js/} : les scripts JavaScript (un par page) ;
  \item \texttt{pages/} : les pages HTML internes ;
  \item \texttt{python/} : les scripts d'IA et le script d'import ;
  \item \texttt{sql/} : les scripts de création et de remplissage de la base.
\end{itemize}

% ======================================================
\section{Conclusion}

% À compléter une fois la fonctionnalité 5 terminée.

\end{document}
